\documentclass[11pt,a4paper]{article}

\usepackage{fontspec}
\setmainfont{DejaVu Sans}
\setsansfont{DejaVu Sans}
\setmonofont{DejaVu Sans Mono}[Scale=0.88]
\defaultfontfeatures{Ligatures=TeX}
\usepackage{microtype}
\usepackage{unicode-math}
\setmathfont{Latin Modern Math}[Scale=1.0]
\usepackage[ngerman]{babel}
\usepackage{geometry}
\geometry{margin=2.5cm}

% --- Zeilenumbruch: sauber wie im Harmonikabuch ---
\sloppy
\emergencystretch=3em
\tolerance=2500
\hbadness=10000
\hfuzz=2pt
\clubpenalty=10000
\widowpenalty=10000

\usepackage{amsmath,amssymb}
\usepackage{parskip}
\usepackage{booktabs}
\usepackage{array}
\usepackage{xcolor}
\usepackage{hyperref}
\hypersetup{
  colorlinks=true,
  linkcolor=blue!50!black,
  urlcolor=blue!50!black,
  pdftitle={FFGFT, das Standardmodell und die Tool-Praxis - eine Abklärung},
  pdfauthor={Johann Pascher}
}

\title{FFGFT, das Standardmodell und die Tool-Praxis\\[2pt]
  \large Eine Abklärung dessen, was FFGFT operativ leistet und nicht leistet}
\author{Johann Pascher\\[2pt]
\small ORCID: \href{https://orcid.org/0009-0000-6518-4064}{0009-0000-6518-4064}\\
\small GitHub: \url{https://github.com/jpascher/T0-Time-Mass-Duality}\\
\small Zenodo (FFGFT v1.0.13): \url{https://zenodo.org/records/20041529}\\
\small DOI: \href{https://doi.org/10.5281/zenodo.20041529}{10.5281/zenodo.20041529}}
\date{Mai 2026}

\begin{document}
\maketitle

\begin{abstract}
\noindent Diese Notiz klärt, was FFGFT im Verhältnis zum operativen Standardmodell-Werkzeugkasten leistet und was nicht. Der Leser soll am Ende drei Dinge auseinanderhalten können: (1) was im Standardmodell-Lagrangian-Formalismus tatsächlich gerechnet wird und mit welchen Werkzeugen, (2) welche Funktion FFGFT als geometrische Reduktion auf einen einzigen dimensionslosen Parameter $\xi = 4/30{,}000$ erfüllt, und (3) wo FFGFT operativ keinen Genauigkeitsvorteil bietet, wo aber konzeptionell und für neue Vorhersagen ein substantieller Beitrag liegt. Die Notiz verfolgt die Doc-206-Disziplin (\S 11/\S 12): klar zu sagen, was die Theorie leistet und was nicht.
\end{abstract}

\section{Anlass: Wer rechnet mit der Lagrangian?}

Die Standardmodell-Lagrangian erscheint in Lehrbüchern in zwei Formen. Die kompakte Hauptformel,
\[
\mathcal{L}_{\mathrm{SM}} = -\tfrac{1}{4}\,F_{\mu\nu}F^{\mu\nu} + i\,\bar\psi\,\!\not{\!D}\psi + \psi_i\,Y_{ij}\,\psi_j\,\phi + \mathrm{h.c.} + |D_\mu\phi|^2 - V(\phi)\,,
\]
ist eine \emph{didaktisch-konzeptionelle} Form, die auf vier Zeilen passt. Die ausgeschriebene Form, in der die drei Eichgruppen $\mathrm{SU}(3)\times\mathrm{SU}(2)\times\mathrm{U}(1)$ komponentenweise expandiert sind, über die drei Generationen summiert ist, mit CKM- und PMNS-Mischungen, allen Yukawa-Kopplungen und allen Lorentz-/Spinor-Indizes, umfasst typischerweise 50--100 Zeilen Formeln mit etwa 19 freien Parametern.

\textbf{Die zentrale Beobachtung:} Niemand rechnet mit der gesamten ausgeschriebenen Lagrangian. Die operative Praxis arbeitet sektorweise und mit einer Industrie von Computer-Werkzeugen, die zwischen Lagrangian-Beschreibung und Vorhersage vermitteln.

\section{Die operative Praxis im Standardmodell}

Die typische Pipeline läuft so:

\begin{enumerate}\itemsep2pt
\item \textbf{Sektorweise Auswahl.} Eine Theoretikerin, die einen LHC-Prozess wie $gg \to H \to \gamma\gamma$ rechnet, nimmt nur die relevanten Teile der Lagrangian: den QCD-Sektor für den Anfangszustand, die Higgs-Yukawa-Kopplung an Top, die elektromagnetische Kopplung für den Endzustand. Die elektroschwache Mischung, die Neutrino-Sektoren, die Lepton-Yukawas usw.\ kommen nicht vor.

\item \textbf{Übersetzung in Feynman-Regeln.} Lagrangian-Terme werden mechanisch in Vertexfaktoren und Propagatoren umgesetzt. Ab hier wird mit Diagrammen gerechnet, nicht mit Lagrangian-Termen.

\item \textbf{Computer-Werkzeuge.} Eine etablierte Toolchain übernimmt die Übersetzung und Rechnung:
\begin{itemize}\itemsep1pt
\item \textbf{FeynRules} -- generiert Feynman-Regeln aus einer Lagrangian-Modelldatei
\item \textbf{FormCalc / FeynCalc} -- Loop-Integrale und Renormierung
\item \textbf{MadGraph / Sherpa / Pythia} -- Event-Generatoren für LHC-Daten
\item \textbf{Lattice-QCD-Codes} -- nicht-perturbative QCD-Rechnungen
\item \textbf{CAMB / CLASS} -- kosmologische Code-Pipelines
\item \textbf{AlterBBN / PArthENoPE} -- Big-Bang-Nukleosynthese-Codes
\end{itemize}

\item \textbf{Effektive Theorien.} Für spezifische Energiebereiche werden reduzierte Lagrangians benutzt: HQET (schwere Quarks), ChPT (chirale Störungstheorie), NRQCD, SMEFT. In ihnen sind die nicht-relevanten Freiheitsgrade ausintegriert.

\item \textbf{Konzeptionelle Funktion der Hauptformel.} Sie wird nicht zur Rechnung benutzt, sondern für Symmetrie-Argumente, Anomalie-Diskussionen, Renormierungsgruppen-Argumente und Konsistenz-Prüfungen.
\end{enumerate}

\textbf{Schlüsselbeobachtung:} Die operative Komplexität liegt nicht in der Lagrangian-Form, sondern in den 19+ freien Parametern, die aus dem Experiment justiert werden müssen, und in der Computer-Vermittlung zwischen Lagrangian und Vorhersage.

\section{Was FFGFT leistet}

FFGFT ist eine geometrische Reduktion. Ein einziger dimensionsloser Parameter,
\[
\xi = \frac{4}{30{,}000} = 1{,}333 \times 10^{-4}\,,
\]
zusammen mit der Skalen-Anker-Beziehung $\tilde T \cdot m = 1$, leitet eine ganze Hierarchie physikalischer Skalen und Konstanten ab, die im Standardmodell als unabhängige freie Parameter auftreten:

\begin{table}[h]
\centering\footnotesize
\begin{tabular}{lcc}
\toprule
\textbf{Größe} & \textbf{Standardmodell-Status} & \textbf{FFGFT-Status} \\
\midrule
$\alpha^{-1}$ & freie Eingabe (gemessen) & abgeleitet aus $\xi$ und Geometrie \\
$\sin^2\theta_W$ & freie Eingabe (gemessen) & exakt $3/13$ \\
Lepton-Massen $m_e, m_\mu, m_\tau$ & 3 freie Yukawas & aus $\xi$-Rekursion (Doc 006) \\
Quark-Massen, CKM & freie Eingaben & aus $\xi$-Geometrie \\
Higgs-Masse, vev & 2 freie Eingaben & aus $16\pi^3$ und $\xi$ (Doc 073) \\
Koide-Verhältnis $Q$ & nicht beachtet & exakt $2/3$ \\
$\eta_{\mathrm{Baryon}}$ & freie Eingabe & $\sim 6\times 10^{-10}$ aus $\xi$ \\
$H_0$ & freie kosm.\ Eingabe & $(\pi/2)\xi^{10}\,m_e c^2/\hbar$ (Doc 165) \\
$L_\xi$ (Casimir) & nicht im SM & $\approx 100\,\mu$m aus $\rho_{\mathrm{CMB}}$ (Doc 091) \\
$\Lambda$ (kosmolog.\ Konstante) & freie Eingabe & geometrisch fixiert \\
\bottomrule
\end{tabular}
\caption{Die Reduktion: 19+ freie Parameter im SM werden zu Konsequenzen einer einzigen geometrischen Struktur in FFGFT.}
\end{table}

\section{Was operativ überflüssig wird}

Aus FFGFT-Sicht wird ein erheblicher Teil der ausgeschriebenen SM-Lagrangian-Komplexität \emph{Bookkeeping einer emergenten Struktur}, deren tatsächliche Quelle die $\xi$-Geometrie ist:

\begin{itemize}\itemsep2pt
\item \textbf{Die 19+ freien Parameter} sind keine unabhängigen Größen, sondern $\xi$-Funktionen.
\item \textbf{Die hierarchische Yukawa-Matrix-Form} (warum $m_\tau / m_e \sim 3500$, nicht 1 oder $10^{10}$) ist Konsequenz der $\xi$-Rekursionstiefe.
\item \textbf{Die spezifische Mischungsstruktur in CKM und PMNS} folgt aus der $\xi$-Geometrie der Generationen.
\item \textbf{Der Higgs-Sektor mit seinen $\lambda$-, $\mu^2$-Postulaten} reduziert sich auf $16\pi^3$- und $\xi$-Strukturen.
\item \textbf{Die spezifische Wahl von $\sin^2\theta_W = 0{,}231$} ist nicht freier Parameter, sondern $3/13$ aus Z$_3$-Geometrie.
\end{itemize}

\section{Was operativ nicht überflüssig wird}

Die SM-Lagrangian-Form als \emph{operationales Werkzeug} bleibt nützlich. Was nicht entfällt:

\begin{itemize}\itemsep2pt
\item \textbf{Feynman-Regeln} -- als operationale Brücke zu experimentellen Wirkungsquerschnitten.
\item \textbf{Renormierungs-Maschinerie} -- für Strahlungskorrekturen.
\item \textbf{Eichinvarianz-Struktur} -- Konsequenz der $\xi$-Geometrie, aber operative Form, in der die Konsequenz manifest wird.
\item \textbf{Die Toolchain} (FeynRules, MadGraph, CAMB usw.) -- für Vergleiche mit Daten.
\end{itemize}

\textbf{Paradigmatischer Wechsel, kein Widerspruch:} Die SM-Lagrangian ist nicht falsch, sondern eine korrekte Beschreibung auf einer mittleren Beschreibungsebene. FFGFT liefert die geometrische Quelle, aus der diese effektive Beschreibung folgt.

\section{Die Genauigkeitshierarchie -- ein zentraler Punkt}

Hier liegt eine Differenzierung, die oft übersehen wird. \textbf{Auf operativer Ebene liefert FFGFT auf einzelne Konstanten keine Genauigkeitsverbesserung gegenüber der direkten Messung.} Beispiel $\alpha^{-1}$:

\begin{align*}
\text{FFGFT-Ableitung:} \quad \alpha^{-1} &= 137{,}035 \quad \text{(4--5 signifikante Stellen)} \\
\text{PDG-Messwert:} \quad \alpha^{-1} &= 137{,}035\,999\,084 \pm 0{,}000\,000\,021 \quad \text{(10+ signifikante Stellen)}
\end{align*}

Für eine Präzisions-QED-Rechnung in MadGraph oder FormCalc will man den PDG-Wert. Eine FFGFT-basierte Eingabe wäre systematisch \emph{ungenauer} als die direkte experimentelle Bestimmung. Das gilt für alle Konstanten, die experimentell auf $\geq 6$ signifikante Stellen bekannt sind: $m_e, m_\mu, m_W, m_Z, \sin^2\theta_W$ usw.

\textbf{Praktische Konsequenz:} Wer FFGFT in die Standard-Toolchain einbauen will, sollte die operativen Eingabewerte aus PDG nehmen und die FFGFT-Ableitungen als \emph{Konsistenz-Check} verwenden. Die Frage ist nicht: ,,Liefert FFGFT bessere LHC-Vorhersagen?'' Die Frage ist: ,,Sind die FFGFT-Ableitungen mit den PDG-Werten innerhalb der Theorie-Genauigkeit kompatibel?''

\section{Wo FFGFT operativ Mehrwert liefert}

Vier Funktionen, in denen FFGFT operativ einen Beitrag leistet, der durch PDG-Messung nicht ersetzbar ist:

\begin{enumerate}\itemsep2pt
\item \textbf{Theoretische Erklärung der Konstanten.} Warum hat $\alpha^{-1}$ den Wert 137,035 und nicht 142,5 oder 119,3? PDG misst es; das SM kann es nicht erklären; FFGFT liefert die geometrische Struktur, aus der der Wert folgt.

\item \textbf{Vorhersagen jenseits der Präzisions-Grenze.} Wo Messwerte ungenau sind, kann FFGFT eine Aussage liefern, die präziser ist als der aktuelle experimentelle Stand. Beispiel: die CHSH-Bell-Abweichung $\sim 10^{-4}$ für massive Teilchen (Doc 023) -- aktuelle Bell-Experimente mit massiven Teilchen haben hier weniger Genauigkeit als die FFGFT-Vorhersage.

\item \textbf{Vorhersagen, die im SM gar nicht existieren.}
\begin{itemize}\itemsep1pt
\item Casimir-Vakuumskala $L_\xi \approx 100\,\mu$m direkt aus $\rho_{\mathrm{CMB}}$ und $\xi$ (Doc 091)
\item Hubble-Konstante $H_0$ aus $\xi^{10}$ (Doc 165), ohne kosmologische Pipeline
\item Spite-Plateau-Erklärungs-Mechanismus über stationäres Reaktionsnetzwerk im $\xi$-Vakuum (Lithium-Briefing)
\item Cosmological constant $\Lambda$ als geometrische Konsequenz, nicht als 120-Größenordnungs-Anomalie
\end{itemize}

\item \textbf{Strukturelle Verbindungen.} FFGFT prädiziert Beziehungen \emph{zwischen} Größen, die im SM als unabhängig behandelt werden:
\[
L_\xi \cdot \frac{H_0}{c} = \frac{\pi}{2}\left(\frac{15}{\pi^2}\right)^{1/4}\frac{m_e c^2}{k_B T_{\mathrm{CMB}}}\,\xi^{41/4}\,, \qquad \frac{41}{4} = 10 + \frac{1}{4}\,.
\]
Eine genauere PDG-Messung der Einzelgrößen $L_\xi, H_0, m_e, T_{\mathrm{CMB}}$ ersetzt diese Vorhersage \emph{zwischen} ihnen nicht.
\end{enumerate}

\section{Tool-Implementation: praktische Andockstellen}

Wenn jemand mit Tool-Zugang FFGFT prüfen oder anwenden will, sind folgende Andockstellen direkt zugänglich:

\begin{enumerate}\itemsep2pt
\item \textbf{FeynRules-Modelldatei für FFGFT.} Die FFGFT-Lagrangian-Form (Doc 019, 067) als \texttt{FFGFT.fr}-Datei. Damit ist FFGFT für alle nachgelagerten Tools (MadGraph, FormCalc usw.) zugänglich. Die operativen Werte sollten aus PDG kommen, die FFGFT-Werte als Konsistenz-Check.

\item \textbf{CAMB / CLASS mit FFGFT-Kosmologie.} $H_0 \approx 66{,}5$ km/s/Mpc, $T_{\mathrm{CMB}}, \eta, \rho_\Lambda$ als Inputs setzen, Vergleich mit Planck-Spektren. Test der Konsistenz der FFGFT-Ableitungen mit CMB-Beobachtungen.

\item \textbf{AlterBBN / PArthENoPE mit Stationaritäts-Annahme.} Statt der zeitlich aufgelösten Freeze-out-Rechnung eine stationäre Netzwerk-Berechnung im $\xi$-Vakuum. Ziel: zeigen, ob das $^7$Li-Plateau natürlich als Attraktor herauskommt. Dies ist genau die offene Arbeit aus dem Lithium-Briefing \S 4.

\item \textbf{Casimir-Setup mit $L_\xi$-Vorhersage.} Direkter experimenteller Test, ob bei Plattenabstand $d \approx L_\xi \approx 100\,\mu$m eine Charakteristik der Casimir-Energiedichte sichtbar ist, die mit $\rho_{\mathrm{CMB}}$ übereinstimmt.

\item \textbf{Bell-Experimente mit massiven Teilchen.} Test der CHSH-Abweichung $\sim 10^{-4}$, etwa mit $^4$He$^*$-Atom-Korrelationen (ANU-Daten).
\end{enumerate}

Diese Aufgaben sind klar definiert. Sie sind \emph{nicht} Aufgaben des FFGFT-Theoretikers, sondern Andockstellen für Implementierer mit Tool-Zugang. FFGFT als geometrische Theorie ist mit der Bereitstellung dieser Andockstellen vollständig; die operationalen Tests sind eine getrennte Aktivität.

\section{Zusammenfassung -- was FFGFT ist und nicht ist}

Eine ehrliche Liste, im Geist der Doc-206-\S 11/\S 12-Disziplin:

\textbf{FFGFT ist nicht:}
\begin{itemize}\itemsep1pt
\item kein Ersatz für Präzisions-Messungen einzelner Konstanten
\item kein schnelleres oder genaueres Rechen-Werkzeug für LHC-Streuamplituden
\item keine Konkurrenz zu MadGraph, FeynRules oder CAMB
\item keine Behauptung, dass die SM-Lagrangian falsch sei
\item keine fertig ausgearbeitete Theorie der nicht-perturbativen QCD oder der Hadronen-Massen
\end{itemize}

\textbf{FFGFT ist:}
\begin{itemize}\itemsep1pt
\item eine Reduktion der freien Parameter von 19+ auf 1
\item eine geometrische Erklärung, warum die Konstanten ihre konkreten Werte haben
\item eine Quelle struktureller Verbindungen zwischen Größen, die im SM unabhängig erscheinen
\item eine Quelle neuer Vorhersagen (Casimir, Bell, Hubble, BBN-Mechanismus) jenseits des SM-Bereichs
\item eine konzeptionelle Klärung dessen, was die SM-Lagrangian beschreibt: eine emergente Struktur, deren Quelle die $\xi$-Geometrie ist
\end{itemize}

\textbf{Die operative Arbeitsteilung:}
\begin{itemize}\itemsep1pt
\item PDG-Messwerte liefern operative Eingaben für Hochpräzisions-Rechnungen.
\item FFGFT liefert konzeptionelle Reduktion und Vorhersagen außerhalb der PDG-Tabelle.
\item Tool-Implementation ist eine eigenständige Aktivität für Praktiker mit Tool-Zugang -- keine Aufgabe des theoretischen Programms.
\end{itemize}

\section*{Quelldokumente im FFGFT-Korpus}

\noindent\small
\begin{tabular}{ll}
\textbf{Doc 003} & Grundlagen: $\xi$-Geometrie und fraktale Raumzeitstruktur \\
\textbf{Doc 006} & Teilchenmassen aus $\xi$-Rekursion \\
\textbf{Doc 019, 067} & Lagrangian-Formulierungen in FFGFT \\
\textbf{Doc 023} & Bell-CHSH-Abweichung für massive Teilchen \\
\textbf{Doc 073} & Higgs-Sektor und $16\pi^3$-Struktur \\
\textbf{Doc 091} & Casimir-CMB-Beziehung, $L_\xi \approx 100\,\mu$m \\
\textbf{Doc 165} & Hubble-Verhältnis $\hbar H_0/(m_e c^2) = (\pi/2)\xi^{10}$ \\
\textbf{Doc 166} & Casimir-Hubble-Brücke, $\xi$-Potenzleiter \\
\textbf{Doc 206} & Triangle-Matrix-Reduktionstheorem \\
\end{tabular}

\medskip
\noindent\small Alle Dokumente verfügbar unter \url{https://github.com/jpascher/T0-Time-Mass-Duality}\\
und als Zenodo-Release v1.0.13: \url{https://zenodo.org/records/20041529} (DOI 10.5281/zenodo.20041529).

\end{document}
